\subsection{Sprint 1 --- 25/08/2023 até 15/09/2023}

Após a retrospectiva da Sprint 0, a equipe recebeu autorização para iniciar uma etapa mais voltada ao desenvolvimento. As definições realizadas anteriormente passaram a servir como base para a construção do software, e as sprints seguintes passaram a exigir maior envolvimento técnico e maior integração entre os membros.

Para esta sprint, foram selecionadas duas \textit{user stories} principais, relacionadas à criação de sessões de Planning Poker e à possibilidade de participar dessas sessões. Com o uso de tecnologias como Node.js e React, essas tarefas pareciam viáveis, mas ainda exigiam adaptação, especialmente para quem estava começando a desenvolver aplicações web.

Inicialmente, minha intenção era contribuir com o backend do projeto. No entanto, por causa da minha inexperiência com Node.js, optei por concentrar meus esforços no frontend, trabalhando na criação de telas para a parte web da aplicação. Fui designado, junto com outros integrantes da AGES I, para desenvolver o menu lateral do projeto.

Essa mudança de foco acabou sendo positiva. A tarefa do menu lateral me proporcionou uma oportunidade concreta de aprendizado em React e Tailwind CSS. Durante a sprint, consegui compreender melhor a organização de componentes, a estilização com classes utilitárias e a importância de manter coerência visual entre as telas. Consegui entregar essa tarefa antes do prazo e, no tempo restante, precisei equilibrar a dedicação à AGES com outras disciplinas da faculdade.

O encerramento da sprint contou com o segundo encontro com os stakeholders, dessa vez sem a presença do Cassio. A reunião foi bastante produtiva e os clientes demonstraram satisfação com as entregas realizadas. Houve pequenas discordâncias relacionadas ao design, mas elas foram resolvidas sem grandes impactos no andamento do projeto.

Ao final da sprint, celebramos um resultado muito positivo. Além de cumprir as metas estabelecidas, a equipe conseguiu avançar em outras frentes, incluindo telas voltadas para dispositivos móveis. Para mim, essa sprint marcou um avanço importante na confiança como desenvolvedor. A experiência com React e Tailwind CSS foi desafiadora, mas mostrou que eu era capaz de aprender tecnologias novas e contribuir com uma entrega concreta dentro de um projeto real.

Também compreendi melhor a importância da comunicação com stakeholders. A satisfação demonstrada na reunião final validou o esforço da equipe e reforçou que o desenvolvimento de software não se limita à implementação técnica: é necessário ouvir, ajustar e alinhar constantemente o produto com as expectativas de quem irá utilizá-lo.
